\section{Occurrence Posterior and Habitable-Zone Integration}\label{sec:occurrence}

\subsection{Bryson occurrence function}

We adopt Model~1 from \citet{Bryson2021}. The differential number of planets
per star per unit $r$ and $I$ is
\begin{equation}
 \lambda(r,I,T\mid\bm\theta)
 =F_0C_1r^\alpha I^\beta T^\gamma g(T),
 \label{eq:lambda}
\end{equation}
where $\bm\theta=(F_0,\alpha,\beta,\gamma)$ is written in manuscript order,
$T$ is in kelvin, and
\begin{equation}
 g(T)=
 \begin{cases}
 10^{-11.839}T^{3.16}, & T\leq5117\ {\rm K},\\
 10^{-16.769}T^{4.49}, & T>5117\ {\rm K}.
 \end{cases}
 \label{eq:gT}
\end{equation}
$C_1$ normalizes the model over $0.5\leq r\leq2.5$,
$0.2\leq I\leq2.2$, and a uniform temperature measure over 3900--6300 K.
Appendix~\ref{app:analytic} gives the analytic normalization.

The constant- and zero-completeness treatments are distinct likelihood
scenarios inherited from the source analysis. They are not samples from one
mixture distribution and are not assigned model probabilities here. The
source interpretation places the unknown long-period completeness between
these cases, but this statement is conditional on the adopted extrapolation
family rather than an assumption-free physical bound.

\subsection{Corrected catalog-error propagation}\label{subsec:measurement}

For a catalog value $\mu$ with reported lower and upper uncertainties
$\sigma_-$ and $\sigma_+$, the v4 primary analysis draws
\begin{align}
 B&\sim{\rm Bernoulli}(1/2),\qquad U=|Z|,\quad Z\sim\mathcal{N}(0,1),\nonumber\\
 X&=\begin{cases}
 \mu-\sigma_-U,&B=0,\\
 \mu+\sigma_+U,&B=1.
 \end{cases}
 \label{eq:measurement_draw}
\end{align}
This construction places the quoted one-sigma lower and upper excursions at
the 15.9th and 84.1st percentiles when the uncertainties are interpreted as
two-sided Gaussian-equivalent errors. It replaces the source notebook's
signed-normal mixture, which does not preserve the intended association
between direction and asymmetric scale. After perturbing instellation, radius,
and $\Teff$, all source-domain conditions are reapplied:
\begin{equation}
 0.2\leq I\leq2.2,
 \quad0.5\leq r\leq2.5,
 \quad3900\leq T\leq6300\ {\rm K}.
 \label{eq:source_domain}
\end{equation}
The source-faithful legacy mechanism is retained only as a regression and
sensitivity mode.

Each outer realization also applies the catalog reliability mechanism before
fitting. Thus the calculation marginalizes over the implemented reliability
and measurement procedure, but not over unmodeled catalog systematics or
occurrence functional forms.

\subsection{Poisson likelihood and adaptive MCMC}\label{subsec:mcmc}

For retained candidates $x_p=(I_p,r_p,T_p)$, the source inhomogeneous-Poisson
log likelihood can be written schematically as
\begin{equation}
 \log\mathcal{L}(\bm\theta)
 =\sum_p\log\lambda(x_p\mid\bm\theta)
 -\int\lambda(x\mid\bm\theta)C(x)\,\dd x,
 \label{eq:poisson_likelihood}
\end{equation}
where $C(x)$ is the branch-specific summed completeness surface evaluated on
the source grid. We retain the source uniform priors inside its numerical
bounds: $0\leq F_0<50000$, $-5\leq\alpha<5$,
$-5\leq\beta<5$, and $-500\leq\gamma<50$, with zero density outside.

For each completeness branch we generate 400 independently seeded outer
reliability/measurement realizations. Each four-parameter conditional
posterior is sampled with 16 ensemble walkers after 1000 discarded burn-in
steps. Autocorrelation time is checked every 1000 production steps. Acceptance
requires, for every parameter, at least $100\tau$, no more than 5\% relative
change between successive $\tau$ estimates, and two consecutive qualifying
checks. The constant production used a 20,000-step ceiling. Five zero-branch
realizations reached that ceiling without satisfying the unchanged gate, so
the entire zero branch was rerun with identical seeds and criteria and a
30,000-step ceiling; no failed realization was discarded.

Adaptive chains have unequal lengths. The aggregation therefore retains 1024
deterministically spaced post-burn rows from every outer realization, assigning
equal weight to each catalog realization. Outer quantile Monte Carlo error is
estimated by resampling whole realizations as clusters; matched contiguous
chain blocks diagnose residual inner-chain error. The two completeness
branches remain separate through this aggregation and the Galactic propagation.

\subsection{Present-day conservative HZ}

For a climatic boundary, the effective stellar flux is
\begin{align}
 S_{\rm eff}(T)={}&S_{{\rm eff},\odot}+aT'
 +bT'^2+cT'^3+dT'^4,\nonumber\\
 T'={}&T-5780\ {\rm K},
 \label{eq:kopparapu}
\end{align}
using the corrected coefficients of \citet{Kopparapu2013,Kopparapu2013Erratum}. The canonical conservative HZ uses the
$1\,\Mearth$ runaway-greenhouse inner boundary from \citet{Kopparapu2014} and
the maximum-greenhouse outer boundary. ``Conservative'' names this prescribed
boundary pair; it is not a universal lower bound over climate models.

Let $S_{\rm out}(T)$ and $S_{\rm in}(T)$ denote the outer and inner flux
boundaries. The broad HZ and narrow target domains are
\begin{equation}
 \mathcal{D}_{\rm HZ}(T)
 =[0.5,1.5]\times[S_{\rm out}(T),S_{\rm in}(T)],
 \label{eq:d_hz}
\end{equation}
\begin{align}
 \mathcal{D}_{\rm EE}(T)=[0.9,1.1]\times
 \big[&\max\{0.9,S_{\rm out}(T)\},\nonumber\\
      &\min\{1.1,S_{\rm in}(T)\}\big],
 \label{eq:d_ee}
\end{align}
with an empty domain when the upper limit does not exceed the lower. Across
5300--6000 K the outer HZ does not clip $I=0.9$; the runaway-greenhouse
boundary crosses $I=1.1$ at 5727.1 K and clips the inner edge at lower
temperatures.

\begin{figure*}[t]
\centering
\includegraphics[width=0.96\textwidth]{v4_hz_boundaries.pdf}
\caption{Conservative habitable-zone boundaries over the selected stellar
temperature range. The broad shaded region lies between the
maximum-greenhouse and $1\,\Mearth$ runaway-greenhouse flux limits. The darker
region is the resulting climate-clipped intersection with the fixed
$0.9\leq I\leq1.1$ narrow target; its hot-side behavior changes where the
runaway boundary crosses $I=1.1$.}
\label{fig:hzboundaries}
\end{figure*}

\subsection{Posterior Galactic estimator}

For $k\in\{{\rm HZ},{\rm EE}\}$ and posterior draw $s$, the expected planets
per star at temperature $T$ are
\begin{equation}
 f_{k,s}(T)=\iint_{\mathcal{D}_k(T)}
 \lambda(r,I,T\mid\bm\theta_s)\,\dd r\,\dd I.
 \label{eq:fk}
\end{equation}
At radial node $R_j$, the occurrence-weighted surface density is
\begin{equation}
 \Sigma_{k,s}(R_j)=\sum_i H_{ij}w_{ij}f_{k,s}(T_{ij}).
 \label{eq:sigma_k}
\end{equation}
The propagated estimator is
\begin{equation}
 \boxed{
 \Lambda_{k,s}^{(\Delta R)}
 =10^6\Delta R\sum_{j=0}^{4}c_j\,2\pi R_j
 \sum_i H_{ij}w_{ij}f_{k,s}(T_{ij}) .}
 \label{eq:galactic_estimator}
\end{equation}
Posterior summaries are quantiles of $\{\Lambda_{k,s}\}$ after equalizing the
outer-realization mixture. Equation~(\ref{eq:galactic_estimator}) preserves the
temperature dependence of both occurrence and HZ clipping; it is not a single
rate multiplied by a host count.

\begin{figure*}[t]
\centering
\includegraphics[width=0.96\textwidth]{v4_occurrence_vs_teff.pdf}
\caption{Temperature-dependent broad-HZ and narrow-target occurrence implied
by the frozen corrected posterior. Lines are posterior medians and shaded
bands are central 68\% intervals. For plotting only, each branch uses a
deterministic 128-draw subsample from every one of its 400 equally weighted
outer realizations (51,200 draws per branch); the numerical tables and
headline posterior summaries remain those of the full frozen propagation.}
\label{fig:occurrence-teff}
\end{figure*}
